
\subsubsection{6.1 Misalignment Findings}

The Spearman ρ = 0.3662 between task-adapted LoRA attention and LM-trained Locret CIS indicates that approximately 86\% of the variance in LoRA attention priority is not explained by Locret CIS priority. This is a large misalignment under any interpretation of "alignment." The consistency across 100 examples (σ = 0.076, no borderline cases) suggests the finding is not a sampling artifact.

Per-layer analysis shows the misalignment is not localized: no transformer layer shows systematically high alignment. The highest per-layer mean ρ values (~0.52–0.53) appear at layers 9 and 11, while early layers (0–2) show the lowest values (~0.17–0.21). This breadth of misalignment suggests the divergence between task-discriminative and next-token-predictive token priority is a global property of the model, not a layer-specific phenomenon.

The GQA expansion artifact (Section 4.4) is a confound that may cause the reported ρ to underestimate the true KV-head-level alignment. The magnitude of this effect is not quantified in the current experiments and is identified as a robustness check for future work.

\subsubsection{6.2 Mechanism Confirmation and Limitations}

The H-M1 PoC demonstrates that: (a) task classification gradients flow through the soft KV budget mask to Locret retaining head weights, and (b) this gradient flow produces a measurable accuracy improvement (+1.50pp over B1) in a single epoch. The improvement did not reach the pre-registered 2.0pp threshold, and the result is classified PARTIAL.

The key limitation of H-M1 is its reduced training scale. The pre-registered protocol called for 3 seeds, 3–5 epochs, and 2000 samples per task; the executed PoC used 1 seed, 1 epoch, and 500 samples. The approximately 4–6× reduction in gradient steps is the most likely explanation for the 0.50pp shortfall, though this explanation is not empirically confirmed.

H-M1 compares JointLoRA-KV against B1 (frozen Locret), not against B3 (sequential LoRA→Locret). The practically relevant claim — that JointLoRA-KV improves over sequential fine-tuning — has no empirical support from the executed experiments.

The H-M2 stability result is restricted to a PoC model (d=64, 2 layers). The argument for stability at full scale rests on the disjoint parameter structure of LoRA (A/B matrices) and Locret (W₁/W₂ per layer), which holds regardless of model scale. However, this architectural argument is not a substitute for empirical confirmation on LLaMA-3.1-8B.

\subsubsection{6.3 Unexecuted Experiments}

Two of five planned experiments were not run. H-M3 (full-scale JointLoRA-KV vs B3 on LongBench-QA + GLUE, 3 seeds) is the primary test of the main hypothesis and was not executed due to compute constraints; all code is implemented and reported as validator-approved. H-M4 (linear probing of retained KV representations) depends on H-M3 and was likewise not executed.

The paper's primary prediction — ≥3\% improvement over B3 on LongBench-QA — therefore has no supporting experimental evidence.

\subsubsection{6.4 Broader Considerations}

Joint training of PEFT adapters and KV eviction policies is a general approach applicable wherever task-specific fine-tuning is combined with memory-efficient inference. The method requires no base model architecture changes — it modifies the training procedure for existing LoRA-compatible and Locret-compatible systems.

A potential concern not investigated in this work is that jointly training the eviction policy on a narrow task distribution could reduce out-of-distribution robustness: if Locret heads are optimized for MNLI token patterns, their eviction decisions on other tasks may degrade. This tradeoff is not assessed.

All experiments use LLaMA-3.1-8B. Generalization to other architectures (Mistral-7B, Qwen-2.5-7B) has not been investigated.

